\section{Семинар 7. Числа Стирлинга, биномиальные суммы и полиномиальная формула}
\noeqref{eq:surj_pie, eq:stirling_2_def, eq:falling_fact, eq:stirling_basis_2}
\noeqref{eq:stirling_basis_1, eq:binom_sq_sum, eq:binom_m_pow, eq:binom_diag_sum}
\noeqref{eq:binom_odd_row_even, eq:half_sum_even, eq:half_sum_odd}
\noeqref{eq:absorb_id, eq:half_sum_weighted, eq:gen_binom_def, eq:neg_binom_id}
\noeqref{eq:multinom_coeff, eq:multinom_thm}

\begin{remark}[Обзор темы]
    Седьмой семинар посвящён классическим задачам распределения различимых 
    объектов по ящикам без пустых мест, числам Стирлинга первого и второго 
    рода, алгебраическим методам вычисления полных и «половинчатых» сумм 
    биномиальных коэффициентов, а также обобщению биномиальных коэффициентов 
    на вещественные и отрицательные верхние индексы и полиномиальной теореме.
\end{remark}

\subsection{Раскладка шаров без пустых ящиков и числа Стирлинга}

\begin{problem}[\textbf{Различимые шары в различимых ящиках без пустых ящиков}]
    Сколькими способами можно разложить $m$ попарно различимых шаров по $n$ 
    попарно различимым ящикам ($m \ge n$) так, чтобы ни один из ящиков не 
    оказался пустым?
\end{problem}

\begin{remark}[Почему наивный двухэтапный выбор ошибочен]
    Нередко при решении этой задачи делают следующую ошибку: сначала 
    выбирают $n$ шаров из $m$ ($C_m^n$ способов) для обязательного 
    заполнения пустых ящиков, а оставшиеся $m - n$ шаров распределяют 
    по $n$ ящикам (например, используя $\bar{C}_n^{m-n}$ для неразличимых 
    шаров или $n^{m-n}$ для различимых):
    \[
        C_m^n \cdot \bar{C}_n^{m-n} = 
        \frac{m!}{n!(m - n)!} \cdot \frac{(m - 1)!}{(m - n)!(n - 1)!}.
    \]
    Этот подход \textbf{неверен} по двум причинам: во-первых, множитель 
    $\bar{C}_n^{m-n}$ относится к \emph{неразличимым} шарам; во-вторых, даже 
    если заменить его на $n! \cdot n^{m-n}$, одно и то же итоговое 
    распределение будет посчитано многократно (если в ящике оказалось два 
    шара $a$ и $b$, то один раз эта конфигурация учтётся, когда шар $a$ 
    был выбран на первом этапе, а шар $b$ — на втором, и второй раз — 
    наоборот). Для точного подсчёта необходимо применить принцип включений и 
    исключений.
\end{remark}

\begin{proof}[Решение методом включений и исключений]~\\
    \begin{quote}
        Каждое распределение $m$ различимых шаров по $n$ различимым ящикам 
        взаимно однозначно соответствует отображению $f: \{1, \dots, m\} \to 
        \{1, \dots, n\}$, ставящему каждому шару в соответствие номер его 
        ящика. Требование отсутствия пустых ящиков означает, что отображение 
        $f$ является \textbf{сюръективным}.

        Обозначим через $U$ множество всех возможных распределений шаров 
        по ящикам без ограничений; его мощность равна $\abs{U} = n^m$. 
        Для каждого ящика с номером $i \in \{1, 2, \dots, n\}$ введём 
        подмножество $A_i \subseteq U$ всех распределений, при которых 
        $i$-й ящик остался \emph{пустым}. Тогда количество распределений 
        без пустых ящиков равно:
        \[
            N = \abs{U \setminus \bigcup_{i=1}^{n} A_i} = 
            n^m - \abs{\bigcup_{i=1}^{n} A_i} = 
            n^m - \sum_{i} \abs{A_i} + \sum_{i < j} \abs{A_i \cap A_j} - \dots
        \]
        Найдём мощность пересечения любых $k$ множеств 
        $A_{j_1} \cap A_{j_2} \cap \dots \cap A_{j_k}$ ($1 \le k \le n$). 
        Условие пустоты выбранных $k$ ящиков означает, что все $m$ шаров 
        произвольным образом раскладываются по оставшимся $n - k$ ящикам:
        \[
            \abs{\bigcap_{s=1}^{k} A_{j_s}} = (n - k)^m.
        \]
        Выбрать $k$ пустых ящиков из $n$ можно $C_n^k$ способами. По 
        формуле включений и исключений получаем:
        \begin{equation} \label{eq:surj_pie}
            N = \sum_{i=0}^{n} (-1)^i C_n^i (n - i)^m. \tag{1}
        \end{equation}
    \end{quote}
\end{proof}

\subsubsection{Числа Стирлинга 2-го и 1-го рода и базисы многочленов}

\begin{definition}
    Пусть теперь $n$ ящиков \textbf{неразличимы} между собой (порядок 
    ящиков не важен). Количество способов разложить $m$ различимых шаров по 
    $n$ неразличимым ящикам так, чтобы ни один ящик не был пуст (то есть 
    количество разбиений $m$-элементного множества на $n$ непустых 
    неупорядоченных подмножеств), называется \textdef{числом Стирлинга 
    2-го рода} и обозначается $S(m, n)$ или $\left\{ \begin{matrix} m \\ n 
    \end{matrix} \right\}$.
\end{definition}

Поскольку все $n$ непустых подмножеств (ящиков) состоят из различных 
элементов, они попарно различны по составу. Следовательно, отождествление 
перестановок $n$ ящиков уменьшает количество вариантов из формулы 
\eqref{eq:surj_pie} ровно в $n!$ раз:
\begin{equation} \label{eq:stirling_2_def}
    S(m, n) = \frac{1}{n!} \sum_{i=0}^{n} (-1)^i C_n^i (n - i)^m. \tag{2}
\end{equation}

\begin{definition}
    \textdef{Убывающей факториальной степенью} (или \textdef{нижним 
    факториалом}) переменной $x$ порядка $k \in \N \cup \{0\}$ называется 
    многочлен степени $k$:
    \begin{equation} \label{eq:falling_fact}
        x^{\underline{k}} = x(x - 1)(x - 2) \dots (x - k + 1) = 
        \frac{x!}{(x - k)!}, \qquad x^{\underline{0}} = 1. \tag{3}
    \end{equation}
    \emph{(В ряде источников сдвиг индекса записывают до множителя $(x - k)$, 
    тогда $x(x - 1)\dots(x - k) = \frac{x!}{(x - k - 1)!} = x^{\underline{k+1}}$).}
\end{definition}

\begin{theorem}[\textbf{Связь степенного и факториального базисов}]
    Числа Стирлинга 2-го рода являются коэффициентами разложения обычной 
    степени $x^n$ по базису убывающих факториальных степеней 
    $x^{\underline{k}}$:
    \begin{equation} \label{eq:stirling_basis_2}
        x^n = \sum_{k=0}^{n} S(n, k) \, x^{\underline{k}}. \tag{4}
    \end{equation}
    Коэффициенты обратного разложения убывающей факториальной степени 
    $x^{\underline{n}}$ по обычным степеням $x^k$ называются 
    \textdef{числами Стирлинга 1-го рода} (обозначаются $s(n, k)$ или 
    $c(k, n)$):
    \begin{equation} \label{eq:stirling_basis_1}
        x^{\underline{n}} = \sum_{k=0}^{n} s(n, k) \, x^k. \tag{5}
    \end{equation}
\end{theorem}

\begin{remark}[\textbf{Матричная интерпретация $P = S \cdot F$}]
    В $(n + 1)$-мерном линейном пространстве многочленов степени не выше $n$ 
    рассмотрим два базиса: стандартный степенной базис 
    $P = (1, x, x^2, \dots, x^n)^T$ и базис убывающих факториальных 
    степеней $F = (1, x^{\underline{1}}, x^{\underline{2}}, \dots, 
    x^{\underline{n}})^T$. Тогда соотношения \eqref{eq:stirling_basis_2} и 
    \eqref{eq:stirling_basis_1} записываются в компактном матричном виде:
    \[
        P = S \cdot F \iff F = S^{-1} \cdot P,
    \]
    где $S = \bigl(S(i, j)\bigr)_{i,j=0}^n$ --- нижнетреугольная матрица 
    чисел Стирлинга 2-го рода с единицами на главной диагонали 
    ($\det S = 1$), а обратная к ней матрица $S^{-1} = \bigl(s(i, j)\bigr)_{i,j=0}^n$ 
    состоит из знаковых чисел Стирлинга 1-го рода.
\end{remark}

\subsection{Приложения формулы бинома Ньютона}

\begin{example}[\textbf{Степени числа $11$ и треугольник Паскаля}] ~
    Представим число $11$ в виде бинома $11 = 10 + 1$ и разложим его 
    натуральные степени по формуле бинома Ньютона:
    \[
        11^n = (10 + 1)^n = \sum_{i=0}^{n} C_n^i 10^i.
    \]
    При $n \le 4$ все биномиальные коэффициенты $C_n^i$ строго меньше $10$ 
    (являются однозначными цифрами), поэтому десятичная запись числа $11^n$ 
    без переносов разрядов буквально воспроизводит $n$-ю строку 
    треугольника Паскаля:
    \begin{align*}
        11^2 &= (10 + 1)^2 = 1 \cdot 10^2 + 2 \cdot 10^1 + 1 \cdot 10^0 = 121, \\
        11^3 &= (10 + 1)^3 = 1 \cdot 10^3 + 3 \cdot 10^2 + 3 \cdot 10^1 + 1 = 1331, \\
        11^4 &= (10 + 1)^4 = 1 \cdot 10^4 + 4 \cdot 10^3 + 6 \cdot 10^2 + 
        4 \cdot 10^1 + 1 = 14641.
    \end{align*}
    Начиная с $n = 5$ ($C_5^2 = C_5^3 = 10$), возникает перенос десятков в 
    старшие разряды: $11^5 = 161051$.
\end{example}

\begin{problem}[\textbf{Рациональные члены в разложении иррационального бинома}]
    Записать биномиальное разложение для выражения 
    $\bigl(\sqrt{2} + \sqrt[3]{3}\bigr)^{100}$ и определить, сколько 
    рациональных (целых) слагаемых содержится в этой сумме после раскрытия 
    скобок.
\end{problem}

\begin{proof}[Решение]~\\
    \begin{quote}
        По формуле бинома Ньютона имеем:
        \[
            \bigl(\sqrt{2} + \sqrt[3]{3}\bigr)^{100} = 
            \sum_{k=0}^{100} C_{100}^k \bigl(\sqrt{2}\bigr)^k 
            \bigl(\sqrt[3]{3}\bigr)^{100-k} = 
            \sum_{k=0}^{100} C_{100}^k \cdot 2^{\frac{k}{2}} \cdot 
            3^{\frac{100-k}{3}}.
        \]
        Поскольку $C_{100}^k \in \N$, а числа $2$ и $3$ взаимно просты, 
        общий член суммы является рациональным (и одновременно целым) 
        числом тогда и только тогда, когда оба показателя степени 
        $\frac{k}{2}$ и $\frac{100 - k}{3}$ являются целыми числами:
        \[
            \begin{cases}
                k \vdots 2, \\
                (100 - k) \vdots 3
            \end{cases}
            \iff
            \begin{cases}
                k \equiv 0 \pmod{2}, \\
                k \equiv 100 \equiv 1 \pmod{3}.
            \end{cases}
        \]
        По китайской теореме об остатках эта система сравнений по взаимно 
        простым модулям $2$ и $3$ равносильна одному сравнению по модулю $6$:
        \[
            k \equiv 4 \pmod{6} \implies k = 6m + 4, \qquad m \in \Z.
        \]
        Учитывая границы индекса суммирования $0 \le k \le 100$:
        \[
            0 \le 6m + 4 \le 100 \iff 0 \le 6m \le 96 \iff 0 \le m \le 16.
        \]
        Переменная $m$ принимает $16 - 0 + 1 = 17$ целых значений 
        ($k = 4, 10, 16, \dots, 100$).

        \textbf{Ответ:} В разложении содержится $17$ рациональных слагаемых.
    \end{quote}
\end{proof}

\subsection{Доказательство биномиальных тождеств и вычисление сумм}

\begin{problem}
    Доказать следующие три классических тождества для биномиальных 
    коэффициентов:
    \begin{align}
        \text{1) } & \sum_{r=0}^{n} (C_n^r)^2 = C_{2n}^n; 
        \label{eq:binom_sq_sum} \tag{6} \\
        \text{2) } & \sum_{k=0}^{n} C_n^k (m - 1)^{n-k} = m^n; 
        \label{eq:binom_m_pow} \tag{7} \\
        \text{3) } & \sum_{j=0}^{k} C_{n-j}^{k-j} = C_{n+1}^k. 
        \label{eq:binom_diag_sum} \tag{8}
    \end{align}
\end{problem}

\begin{proof}[Решение]~\\
    \begin{quote}
        \textbf{1. Доказательство тождества \eqref{eq:binom_sq_sum}:}\\
        Прямое раскрытие суммы квадратов факториалов 
        $\sum_{r=0}^n \left(\frac{n!}{r!(n - r)!}\right)^2$ технически 
        громоздко. Вместо этого рассмотрим алгебраическое тождество для 
        многочленов от переменной $x$:
        \[
            (1 + x)^{2n} = (1 + x)^n \cdot (1 + x)^n = 
            \left(\sum_{r=0}^{n} C_n^r x^r\right)^2 = 
            \left(\sum_{r=0}^{n} C_n^r x^r\right) 
            \left(\sum_{s=0}^{n} C_n^s x^s\right).
        \]
        Приравняем коэффициенты при степени $x^n$ в левой и правой частях 
        равенства:
        \begin{itemize}
            \item В левой части коэффициент при $x^n$ в разложении 
            $(1 + x)^{2n}$ равен $C_{2n}^n$.
            \item В правой части при перемножении двух скобок моном 
            $x^r \cdot x^s = x^n$ получается при $s = n - r$ 
            ($r = 0, 1, \dots, n$). Следовательно, коэффициент при $x^n$ 
            равен свёртке $\sum_{r=0}^n C_n^r C_n^{n-r}$.
        \end{itemize}
        Используя свойство симметрии $C_n^{n-r} = C_n^r$, получаем:
        \[
            C_{2n}^n = \sum_{r=0}^{n} C_n^r \cdot C_n^{n-r} = 
            \sum_{r=0}^{n} (C_n^r)^2.
        \]

        \textbf{2. Доказательство тождества \eqref{eq:binom_m_pow}:}\\
        Запишем формулу бинома Ньютона $(x + y)^n = \sum_{k=0}^n C_n^k x^k y^{n-k}$ 
        и подставим в неё значения $x = 1$ и $y = m - 1$:
        \[
            \sum_{k=0}^{n} C_n^k \cdot 1^k \cdot (m - 1)^{n-k} = 
            \bigl(1 + (m - 1)\bigr)^n = m^n.
        \]

        \textbf{3. Доказательство тождества \eqref{eq:binom_diag_sum}:}\\
        Приведём два способа доказательства (оба намечены в конспекте):
        \begin{itemize}
            \item \emph{Способ А (сведение к суммированию по верхнему 
            индексу):} По свойству симметрии $C_{n-j}^{k-j} = C_{n-j}^{(n-j)-(k-j)} 
            = C_{n-j}^{n-k}$. Выполним замену индекса суммирования 
            $i = n - j$ (когда $j$ меняется от $0$ до $k$, индекс $i$ 
            пробегает значения от $n - k$ до $n$):
            \[
                \sum_{j=0}^{k} C_{n-j}^{k-j} = 
                \sum_{i=n-k}^{n} C_i^{n-k} = C_{n+1}^{n-k+1} = C_{n+1}^k,
            \]
            где мы воспользовались свойством 9 из Лекции 4 (правилом 
            «хоккейной клюшки») и симметрией биномиальных коэффициентов.
            \item \emph{Способ Б (телескопическое суммирование по правилу 
            Паскаля):} Из правила Паскаля выразим $C_{n-j}^{k-j} = 
            C_{n-j+1}^{k-j} - C_{n-j}^{k-j-1}$ (при $j < k$, а при $j = k$ 
            имеем $C_{n-k}^0 = C_{n-k+1}^0 = 1$). Складывая эти равенства 
            по $j = 0, \dots, k$, видим, что все промежуточные слагаемые 
            взаимно уничтожаются, оставляя $C_{n+1}^k$.
        \end{itemize}
    \end{quote}
\end{proof}

\begin{problem}
    Доказать тождество для суммы чётных биномиальных коэффициентов 
    нечётной строки $2n + 1$:
    \begin{equation} \label{eq:binom_odd_row_even}
        \sum_{k=0}^{n} C_{2n+1}^{2k} = 4^n \tag{9}
    \end{equation}
\end{problem}

\begin{proof}[Решение]~\\
    \begin{quote}
        Применим к каждому слагаемому $C_{2n+1}^{2k}$ (при $k \ge 1$) 
        правило Паскаля:
        \[
            C_{2n+1}^{2k} = C_{2n}^{2k} + C_{2n}^{2k-1},
        \]
        а при $k = 0$ заметим, что $C_{2n+1}^0 = 1 = C_{2n}^0$ (можно 
        также формально считать $C_{2n}^{-1} = 0$). Подставляя это 
        разложение в сумму \eqref{eq:binom_odd_row_even}, получаем сумму 
        \emph{всех} подряд идущих биномиальных коэффициентов $2n$-й строки 
        (как с чётными, так и с нечётными нижними индексами):
        \[
            \sum_{k=0}^{n} C_{2n+1}^{2k} = 
            C_{2n}^0 + \sum_{k=1}^{n} \bigl(C_{2n}^{2k-1} + C_{2n}^{2k}\bigr) = 
            \sum_{j=0}^{2n} C_{2n}^j = (1 + 1)^{2n} = 2^{2n} = 4^n.
        \]
    \end{quote}
\end{proof}

\subsubsection{Вычисление «половинчатых» биномиальных сумм}

\begin{problem}
    Вычислить следующие суммы по первой половине строки треугольника 
    Паскаля:
    \begin{align}
        \text{a) } & \sum_{k=0}^{n} C_{2n}^k = 
        2^{2n-1} + \frac{1}{2} C_{2n}^n = 2^{2n-1} + C_{2n-1}^n; 
        \label{eq:half_sum_even} \tag{10} \\
        \text{b) } & \sum_{k=0}^{n} C_{2n+1}^k = 2^{2n}; 
        \label{eq:half_sum_odd} \tag{11} \\
        \text{c) } & \sum_{k=0}^{n} k \cdot C_{2n}^k = 
        2n \sum_{k=0}^{n-1} C_{2n-1}^k = n \cdot 2^{2n-1}. 
        \label{eq:half_sum_weighted} \tag{12}
    \end{align}
\end{problem}

\begin{proof}[Решение]~\\
    \begin{quote}
        \textbf{Пункт a): Вывод формулы \eqref{eq:half_sum_even}.}\\
        Запишем полную сумму коэффициентов чётной строки $2n$, равную 
        $2^{2n}$, и выделим в ней центральный биномиальный коэффициент 
        $C_{2n}^n$:
        \[
            2^{2n} = \sum_{k=0}^{2n} C_{2n}^k = 
            \sum_{k=0}^{n-1} C_{2n}^k + C_{2n}^n + \sum_{k=n+1}^{2n} C_{2n}^k.
        \]
        В силу симметрии $C_{2n}^{2n-k} = C_{2n}^k$ сумма первых $n$ 
        слагаемых равна сумме последних $n$ слагаемых:
        \[
            \sum_{k=n+1}^{2n} C_{2n}^k = \sum_{j=0}^{n-1} C_{2n}^{2n-j} = 
            \sum_{j=0}^{n-1} C_{2n}^j.
        \]
        Обозначим искомую сумму от $k = 0$ до $n$ через 
        $S = \sum_{k=0}^n C_{2n}^k$. Тогда $\sum_{k=0}^{n-1} C_{2n}^k = 
        S - C_{2n}^n$, и уравнение для полной суммы принимает вид:
        \[
            (S - C_{2n}^n) + C_{2n}^n + (S - C_{2n}^n) = 2^{2n} \iff 
            2S - C_{2n}^n = 2^{2n} \implies S = 2^{2n-1} + \frac{1}{2} C_{2n}^n.
        \]
        Преобразуем центральный коэффициент при $n \ge 1$:
        \[
            C_{2n}^n = \frac{(2n)!}{n! \, n!} = 
            \frac{2n \cdot (2n - 1)!}{n \cdot (n - 1)! \, n!} = 
            2 \cdot \frac{(2n - 1)!}{(n - 1)! \, n!} = 2 C_{2n-1}^n.
        \]
        Деля на $2$, получаем $\frac{1}{2} C_{2n}^n = C_{2n-1}^n$, откуда 
        окончательно $S = 2^{2n-1} + C_{2n-1}^n$.

        \textbf{Пункт b): Вывод формулы \eqref{eq:half_sum_odd} (два 
        способа).}
        \begin{itemize}
            \item \emph{Способ 1 (через симметрию нечётной строки):} 
            В строке с нечётным номером $2n + 1$ всего $2n + 2$ элемента 
            (центрального одиночного элемента нет). По свойству симметрии 
            $C_{2n+1}^k = C_{2n+1}^{2n+1-k}$ сумма первой половины строки 
            ($k = 0, \dots, n$) равна сумме второй половины строки 
            ($k = n + 1, \dots, 2n + 1$):
            \[
                2 \sum_{k=0}^{n} C_{2n+1}^k = \sum_{k=0}^{2n+1} C_{2n+1}^k = 
                2^{2n+1} \implies \sum_{k=0}^{n} C_{2n+1}^k = 2^{2n}.
            \]
            \item \emph{Способ 2 (через правило Паскаля и чётные индексы, 
            как в конспекте):} 
            По правилу Паскаля $\sum_{k=0}^n C_{2n+1}^k = 
            \sum_{k=0}^n \bigl(C_{2n}^k + C_{2n}^{k-1}\bigr)$. С другой 
            стороны, заменяя в этой сумме $C_{2n}^k = C_{2n}^{2n-k}$ и 
            $C_{2n}^{k-1} = C_{2n}^{2n-k+1}$, заметим, что каждая пара 
            сворачивается в элемент $(2n + 1)$-й строки, и сумма 
            $\sum_{k=0}^n C_{2n+1}^k$ совпадает с суммой по всем чётным 
            индексам $\sum_{k=0}^n C_{2n+1}^{2k}$, которая по доказанному 
            в \eqref{eq:binom_odd_row_even} равна $4^n = 2^{2n}$.
        \end{itemize}

        \textbf{Пункт c): Вывод формулы \eqref{eq:half_sum_weighted}.}\\
        Докажем предварительно важное \textdef{тождество поглощения} для 
        $k \ge 1$:
        \begin{equation} \label{eq:absorb_id}
            k \cdot C_N^k = \frac{k \cdot N!}{k!(N - k)!} = 
            N \cdot \frac{(N - 1)!}{(k - 1)!\bigl((N - 1) - (k - 1)\bigr)!} = 
            N \cdot C_{N-1}^{k-1}. \tag{13}
        \end{equation}
        Применим тождество \eqref{eq:absorb_id} при $N = 2n$ к каждому 
        слагаемому искомой суммы (нулевое слагаемое при $k = 0$ равно $0$):
        \[
            \sum_{k=0}^{n} k \cdot C_{2n}^k = 
            \sum_{k=1}^{n} 2n \cdot C_{2n-1}^{k-1} = 
            2n \sum_{j=0}^{n-1} C_{2n-1}^j,
        \]
        где сделана замена индекса $j = k - 1$. Заметим, что сумма 
        $\sum_{j=0}^{n-1} C_{2n-1}^j$ представляет собой сумму первой 
        половины нечётной строки с номером $2(n - 1) + 1 = 2n - 1$. 
        Следовательно, по формуле пункта b) \eqref{eq:half_sum_odd} (с 
        заменой $n$ на $n - 1$) имеем:
        \[
            \sum_{j=0}^{n-1} C_{2n-1}^j = 2^{2(n-1)} = 2^{2n-2}.
        \]
        Подставляя это значение, окончательно находим:
        \[
            \sum_{k=0}^{n} k \cdot C_{2n}^k = 2n \cdot 2^{2n-2} = 
            n \cdot 2^{2n-1}.
        \]
    \end{quote}
\end{proof}

\subsection{Обобщённый биномиальный и полиномиальный коэффициенты}

\subsubsection{Обобщённый биномиальный коэффициент}

\begin{definition}
    Для произвольного вещественного (или комплексного) числа $r \in \R$ и 
    любого целого неотрицательного числа $k \in \N \cup \{0\}$ 
    \textdef{обобщённым биномиальным коэффициентом} называется выражение:
    \begin{equation} \label{eq:gen_binom_def}
        C_r^k = \binom{r}{k} = 
        \frac{r(r - 1)(r - 2) \dots (r - k + 1)}{k!} = 
        \frac{r^{\underline{k}}}{k!}, \qquad C_r^0 = 1. \tag{14}
    \end{equation}
\end{definition}

\begin{properties}[Свойства обобщённых биномиальных коэффициентов] ~
    \begin{enumerate}
        \item \textbf{Обращение верхнего индекса (формула для 
        отрицательного верхнего индекса):} 
        Вынося знак «$-$» из каждого из $k$ сомножителей в числителе 
        определения \eqref{eq:gen_binom_def}, получаем связь с сочетаниями 
        с повторениями:
        \begin{equation} \label{eq:neg_binom_id}
            C_{-r}^k = \frac{(-r)(-r - 1) \dots (-r - k + 1)}{k!} = 
            (-1)^k \frac{r(r + 1)(r + 2) \dots (r + k - 1)}{k!} = 
            (-1)^k C_{r+k-1}^k. \tag{15}
        \end{equation}
        \item \textbf{Сохранение правила Паскаля:} 
        Для любого $r \in \R$ и любого натурального $k \ge 1$ выполняется 
        тождество:
        \[
            C_r^k = C_{r-1}^k + C_{r-1}^{k-1}.
        \]
    \end{enumerate}
\end{properties}

\subsubsection{Полиномиальный коэффициент и полиномиальная теорема}

\begin{definition}
    Пусть $n \in \N \cup \{0\}$, а $k_1, k_2, \dots, k_m$ --- целые 
    неотрицательные числа, сумма которых равна $n$ 
    ($k_1 + k_2 + \dots + k_m = n$). \textdef{Полиномиальным 
    (мультиномиальным) коэффициентом} называется число:
    \begin{equation} \label{eq:multinom_coeff}
        C_n^{k_1, k_2, \dots, k_m} = \binom{n}{k_1, k_2, \dots, k_m} = 
        \frac{n!}{k_1! \, k_2! \, \dots \, k_m!}. \tag{16}
    \end{equation}
\end{definition}

\begin{remark}[Комбинаторный смысл полиномиального коэффициента]
    Полиномиальный коэффициент $C_n^{k_1, k_2, \dots, k_m}$ равен:
    \begin{itemize}
        \item числу \textbf{перестановок с повторениями} длины $n$, 
        составленных из элементов $m$ различных типов, где элемент $i$-го 
        типа входит ровно $k_i$ раз;
        \item числу способов разложить $n$ различимых предметов по $m$ 
        различимым ящикам так, чтобы в $1$-м ящике оказалось ровно $k_1$ 
        предметов, во $2$-м --- ровно $k_2$ предметов, $\dots$, в $m$-м --- 
        ровно $k_m$ предметов.
    \end{itemize}
\end{remark}

\begin{theorem}[\textbf{Полиномиальная формула}]
    Для любых чисел $a_1, a_2, \dots, a_m$ и любого целого неотрицательного 
    показателя $n \in \N \cup \{0\}$ справедливо разложение:
    \begin{equation} \label{eq:multinom_thm}
        (a_1 + a_2 + \dots + a_m)^n = 
        \sum_{\substack{k_1 + k_2 + \dots + k_m = n \\ k_i \ge 0}} 
        C_n^{k_1, k_2, \dots, k_m} \cdot 
        a_1^{k_1} \cdot a_2^{k_2} \cdot \dots \cdot a_m^{k_m}. \tag{17}
    \end{equation}
\end{theorem}

\begin{proof}[Доказательство]~\\
    \begin{quote}
        Представим степень $(a_1 + \dots + a_m)^n$ в виде произведения $n$ 
        одинаковых скобок. При раскрытии скобок каждое слагаемое вида 
        $a_1^{k_1} a_2^{k_2} \dots a_m^{k_m}$ (где $k_1 + \dots + k_m = n$) 
        получается, когда переменная $a_1$ выбирается ровно из $k_1$ скобок, 
        переменная $a_2$ --- из $k_2$ скобок из числа оставшихся, и так 
        далее, вплоть до переменной $a_m$, выбираемой из последних $k_m$ 
        скобок. По правилу умножения количество способов такого выбора равно:
        \[
            C_n^{k_1} \cdot C_{n-k_1}^{k_2} \cdot \dots \cdot C_{k_m}^{k_m} = 
            \frac{n!}{k_1! \, \cancel{(n - k_1)!}} \cdot 
            \frac{\cancel{(n - k_1)!}}{k_2! \, \cancel{(n - k_1 - k_2)!}} 
            \cdot \dots \cdot \frac{\cancel{k_m!}}{k_m! \, 0!} = 
            \frac{n!}{k_1! \, k_2! \, \dots \, k_m!} = 
            C_n^{k_1, k_2, \dots, k_m},
        \]
        что доказывает формулу \eqref{eq:multinom_thm}.
    \end{quote}
\end{proof}